\chapter{Scherzugversuch}

Im Rahmen der von der Siemens AG beauftragten Versuchsreihe wurde der Scherzugversuch als statische Voruntersuchung für die nachfolgenden Wöhlerversuche durchgeführt.
Ziel ist nicht die Bestimmung einer werkstoffspezifischen Zugfestigkeit, sondern die Ermittlung der maximal übertragbaren
Kraft (hier der vernieteten Verbindung). Wie in \Chapref{chap:Theorie_Scherzugversuch} erläutert, dient die im Scherzugversuch
ermittelte statische Tragfähigkeit als Orientierungsgröße für die
Festlegung der Lastniveaus des Wöhlerversuchs.
 Diese statische Tragfähigkeit dient als Orientierungsgröße für die Festlegung der Lastniveaus im Wöhlerversuch, die unterhalb der statischen Versagenslast liegen müssen. Zusätzlich wird die auftretende Versagensart dokumentiert. Da der Schwerpunkt der Arbeit auf der standardisierten Aufbereitung und Auswertung der Messdaten der Hydropulsanlage liegt, werden der
Versuchsaufbau und die Durchführung des Scherzugversuchs im Folgenden nur in dem
für die Einordnung und Nachvollziehbarkeit erforderlichen Umfang beschrieben.

	 \section{Versuchsaufbau}
Die Scherzugversuche wurden am Technologie Campus in Neustadt an der Donau durchgeführt. Die dort verwendete Standprüfmaschine der Firma Zwick/Roell ist in \Figref{fig:Zugpruefmaschine} zu sehen. 

    \begin{figure}[H]
        \begin{center}
            \includegraphics[width=0.27\textwidth]{Zugpruefmaschine_Zwick}
            \caption{Zugprüfmaschine Zwick/Roell Z250 (eigene Aufnahme).}
	  \label{fig:Zugpruefmaschine}
        \end{center}
    \end{figure}

Während des Versuchs wurden die aufgebrachte Kraft und der Traversenweg kontinuierlich aufgezeichnet.

	\section{Proben und Vorbereitung}
	Untersucht wurden drei Probenchargen mit jeweils fünf vernieteten Prüfkörpern. Alle
Proben bestanden aus zwei über eine Blindnietverbindung gefügten Blechen. Der
Bohrungsdurchmesser betrug bei sämtlichen Proben \SI{6.4}{\milli\metre}. Die Chargen
unterschieden sich hinsichtlich des Werkstoffs und der Blechdicken. Eine Übersicht ist
in \Tabref{tab:scherzug_probenchargen} angegeben.

\begin{table}[H]
    \centering
    \caption{Übersicht der im Scherzugversuch untersuchten Probenchargen.}
    \label{tab:scherzug_probenchargen}
    \resizebox{\textwidth}{!}{%
        \begin{tabular}{c c c c c}
            \hline
            Charge & Werkstoff & Blechdicke 1 & Blechdicke 2 & Probenanzahl \\
            \hline
            1 & AlMg3 & \SI{3}{\milli\metre} & \SI{3}{\milli\metre} & 5 \\
            2 & Chrom-Nickel-Stahl & \SI{2}{\milli\metre} & \SI{5}{\milli\metre} & 5 \\
            3 & Chrom-Nickel-Stahl & \SI{3}{\milli\metre} & \SI{3}{\milli\metre} & 5 \\
            \hline
        \end{tabular}%
    }
\end{table}

Die Aluminiumproben der ersten Charge sind exemplarisch in \Figref{fig:Zugproben Aluminium} dargestellt.

    \begin{figure}[H]
        \begin{center}
            \includegraphics[width=0.75\textwidth]{Zugproben_Aluminium}
            \caption{Zugproben Aluminium.}
	  \label{fig:Zugproben Aluminium}
        \end{center}
    \end{figure}

Zur eindeutigen Zuordnung der Messdaten, Bruchbilder und Prüfprotokolle wurde für
sämtliche Proben eine einheitliche Probenbezeichnung verwendet. 
\clearpage
Die Bezeichnung \texttt{Z\_P05\_AlMg3\_3\_3mm\_D6.4\_kpl} enthält beispielsweise die Probenart, die Probennummer, den Werkstoff, die beiden Blechdicken, den Bohrungsdurchmesser und den
Bearbeitungsstatus. \\Diese eindeutige Benennung unterstützt die spätere standardisierte
Zuordnung und Weiterverarbeitung der Versuchsdaten.

	 \section{Durchführung}
	\subsection{Versuchsparameter}
	Damit die Nietverbindung möglichst ohne zusätzliche Biegebeanspruchung geprüft wird,
muss die Kraftlinie in der Ebene der Scherfuge liegen. Hierzu wurden Ausgleichsbleche
verwendet, deren Dicke an die jeweilige Probengeometrie angepasst war. Vor der
Versuchsdurchführung wurden die Kontaktflächen gereinigt und die Spannbacken parallel
sowie zentrisch ausgerichtet. Die zentrische und biegemomentarme Einspannung sowie die
gewählte Prüfgeschwindigkeit orientieren sich an den Vorgaben der DIN~EN~ISO~14589
für die mechanische Prüfung von Blindnieten~\cite[Seite~7]{Blindniete_Norm}.

Die wesentlichen Versuchsparameter sind in \Tabref{tab:scherzug_versuchsparameter} zusammengefasst.

\begin{table}[H]
    \centering
    \caption{Wesentliche Parameter der Scherzugversuche.}
    \label{tab:scherzug_versuchsparameter}
    \begin{tabular}{l l}
        \hline
        Prüfparameter & Eingestellter Wert \\
        \hline
        Prüfgeschwindigkeit & \SI{10}{\milli\metre\per\minute} \\
        Freie Prüflänge & \SI{80}{\milli\metre} \\
        Vorkraft & ca. \SI{120}{\newton} \\
        Abbruchkriterium & Versagen der Nietverbindung oder des Blechs \\
        \hline
    \end{tabular}
\end{table}

	\subsection{Versuchsablauf}
	
	Zu Beginn jedes Versuchs wurde die Probe zentrisch in die Spannbacken eingesetzt und
über die angepassten Ausgleichsbleche ausgerichtet. Für die erste Probencharge kamen
zunächst Spannbacken des Typs 8406 zum Einsatz. Da sich einzelne Proben beim Spannen
leicht verzogen, wurden für die beiden nachfolgenden Chargen größere Spannbacken des
Typs 8507 verwendet. Anschließend wurde die freie Prüflänge kontrolliert, die
Schutztür geschlossen und die Vorkraft von etwa \SI{120}{\newton} aufgebracht.

Nach dem Start der Messung wurde die Traverse mit konstanter Geschwindigkeit bewegt,
bis die Nietverbindung oder eines der Bleche versagte. Für jede Probe wurden der
Kraft-Weg-Verlauf, die Maximalkraft und die beobachtete Versagensart dokumentiert.
Diese Größen bilden die Grundlage der folgenden Auswertung.

	 \section{Auswertung}

Die Versuchsergebnisse werden mithilfe von Kraft-Weg-Diagrammen dargestellt. Für jede
Probe wird die Maximalkraft aus dem jeweiligen Kraftverlauf bestimmt. Zur Beschreibung
der Streuung innerhalb einer Charge werden der arithmetische Mittelwert, die
Standardabweichung und der Variationskoeffizient verwendet. Die zugehörigen
statistischen Grundlagen wurden in Kapitel~2 erläutert.

\subsection*{Charge 1: AlMg3, Blechdicken \SI{3}{\milli\metre}/\SI{3}{\milli\metre}}

Die Kraft-Weg-Verläufe der Aluminiumproben sind in
\Figref{fig:Diagramm_AlMg3_3} dargestellt. Für Probe~1 wurde kein verwertbarer
Messdatensatz gespeichert, weshalb die statistische Auswertung dieser Charge auf den
vier vollständig aufgezeichneten Versuchen basiert.
	  
    \begin{figure}[H]
        \begin{center}
            \includegraphics[width=0.8\textwidth]{Scherzugversuch_diagramm_AlMg3_3}
            \caption{Scherzugversuch - Proben 2-5 (AlMg3\_3\_3mm, D6.4).}
	  \label{fig:Diagramm_AlMg3_3}
        \end{center}
    \end{figure}
    
Die mittlere Maximalkraft beträgt 9770\,$N$. Mit einer
Standardabweichung von 40,8\,$N$ und einem Variationskoeffizienten von
0,42\,$\%$ weisen die aufgezeichneten Ergebnisse eine geringe Streuung auf.
Bei den geprüften Aluminiumproben trat das Versagen jeweils im Blech auf.

\subsection*{Charge 2: Chrom-Nickel-Stahl, Blechdicken
\SI{2}{\milli\metre}/\SI{5}{\milli\metre}}

Die Kraft-Weg-Verläufe der zweiten Charge sind in
\Figref{fig:Diagramm_CrNi2_5} abgebildet. Die Proben erreichten das Ausfallkriterium bei einem
geringeren Traversenweg als die Aluminiumproben. Bei allen fünf Proben trat ein
Nietbruch auf. Des Weiteren sind am Anfang der Kurven Zacken zu erkennen, die ihre Ursache vermutlich in einem leichten Verrutschen der Proben haben. 
    
Bei den Chrom-Nickel-Proben mit zwei und fünf Millimeter Blechstärke zeigt sich in \Figref{fig:Diagramm_CrNi2_5}, dass die Proben nach einem deutlich geringeren Weg bereits versagen. Des Weiteren sind am Anfang der Kurven Zacken zu erkennen, die ihre Ursache vermutlich darin haben, dass die Proben leicht verrutscht sind. 
        \begin{figure}[H]
        \begin{center}
            \includegraphics[width=0.8\textwidth]{Scherzugversuch_diagramm_CrNi2_5}
            \caption{Scherzugversuch - Proben 1-5 (CrNi\_2\_5mm, D6.4).}
	  \label{fig:Diagramm_CrNi2_5}
        \end{center}
    \end{figure}
Die mittlere Maximalkraft liegt bei 11000\,$N$, die Standardabweichung bei 328\,$N$ und der Variationskoeffizient bei 2,99\,$\%$. 

\subsection*{Charge 3: Chrom-Nickel-Stahl, Blechdicken
\SI{3}{\milli\metre}/\SI{3}{\milli\metre}}

Die dritte Charge weist die höchsten gemessenen Maximalkräfte auf. Die zugehörigen
Kraft-Weg-Verläufe sind in \Figref{fig:Diagramm_CrNi3_3} dargestellt. Auch bei
dieser Charge versagten sämtliche Proben durch einen Nietbruch.

            \begin{figure}[H]
        \begin{center}
            \includegraphics[width=0.8\textwidth]{Scherzugversuch_diagramm_CrNi3_3}
            \caption{Scherzugversuch - Proben 1-5 (CrNi\_3\_3mm, D6.4).}
	  \label{fig:Diagramm_CrNi3_3}
        \end{center}
    \end{figure}

Die mittlere Maximalkraft beträgt 14900\,$N$, die Standardabweichung 393\,$N$ und der Variationskoeffizient 2,63\,$\%$. 

\subsection*{Zusammenfassung und Einordnung}

\Tabref{tab:scherzug_ergebnisse} fasst die wesentlichen Ergebnisse der drei
Probenchargen zusammen.

\begin{table}[H]
    \centering
    \caption{Zusammenfassung der Ergebnisse der Scherzugversuche.}
    \label{tab:scherzug_ergebnisse}
    \resizebox{\textwidth}{!}{%
        \begin{tabular}{l c c c c}
            \hline
            Probencharge &
            mittlere Maximalkraft &
            Standardabweichung &
            Variationskoeffizient &
            Versagensart \\
            \hline
            AlMg3, 3/3 mm &
            \SI{9770}{\newton} &
            \SI{40.8}{\newton} &
            \SI{0.42}{\percent} &
            Blechversagen \\
            CrNi, 2/5 mm &
            \SI{11000}{\newton} &
            \SI{328}{\newton} &
            \SI{2.99}{\percent} &
            Nietbruch \\
            CrNi, 3/3 mm &
            \SI{14900}{\newton} &
            \SI{393}{\newton} &
            \SI{2.63}{\percent} &
            Nietbruch \\
            \hline
        \end{tabular}%
    }
\end{table}

 Die Ergebnisse zeigen, dass sich sowohl die maximal übertragbare Kraft als auch die
Versagensart in Abhängigkeit von Werkstoff und Blechdicken unterscheiden. Ein
unmittelbarer werkstofflicher Vergleich allein anhand des Traversenwegs ist jedoch nur
eingeschränkt möglich, da sich neben dem Werkstoff auch die Geometrie und die
Versagensart der Chargen unterscheiden.

Für die nachfolgenden Wöhlerversuche ist insbesondere die statische Tragfähigkeit der
jeweils untersuchten Nietverbindung relevant. Die im Scherzugversuch ermittelten
Maximalkräfte dienen daher als obere Orientierungsgröße für die Festlegung der
dynamischen Lastniveaus. Die Lasten des Wöhlerversuchs werden unterhalb dieser
statischen Versagenslasten gewählt, sodass nicht ein unmittelbares statisches Versagen,
sondern das Ermüdungsverhalten der Verbindung untersucht wird. Die vollständigen
Prüfprotokolle der Scherzugversuche sind in \Anhref{anhang:pruefprotokoll} enthalten.
